\section{Introduction}
Consider editing an object's shape from a single target image while retaining
appearance evidence from its original views. The edited target and the
sources depict different geometric states. Even with calibrated cameras, a
source feature projected along its original ray need not describe the desired
edited region. A pretrained multi-view generator supplies a useful 3D prior,
but its native spatial associations do not resolve this intentional conflict.

Semantic correspondence can reconnect related regions across the edit.
However, selecting one match can discard a useful alternative, while mixing
many matches can transfer irrelevant values. A better retrieval score also
does not establish that the generated surface or material improves. The
delivery path matters: adapted evidence enters an existing projection-based
conditioning function whose target contribution and source features already
have a trained numerical scale.
Figure~\ref{fig:teaser} illustrates the distinct target and source evidence
and the evaluated view roles.

For paired supervision, we construct DeformNeXt-10k using donor-guided
GeoStyle deformation with source assembly constraints and final-pair
screening. The dataset retains exact canonical mesh correspondence while
providing 10,165 accepted edits from 3,483 source assets, rather than one
fixed edit per object. Section~\ref{sec:data-construction} describes the
construction, split boundary, and limitations of its shortened production
optimization budget.

We study source adaptation within that native function. Voxels projecting to
the same source-camera patch share one evidence decision. A frozen matcher
constructs semantic proposals from the observed target members of that ray;
the selected SS design aggregates native values without a new value/output
projection. A source-contrast retargeter then supplies a gated correction to
the native carrier. Its zero initialization preserves the original function,
and rays without proposal support retain the pretrained prior. The target
conditioning stream is retained; exact target output preservation remains
an evaluation question.

Our experiments distinguish stage objectives. SS candidates are evaluated
for geometric alignment under common frozen downstream models. PBR
candidates are evaluated for source appearance on identical generated
geometry, guides, noise, and rendering triangles. This produces a measured
composition of Fixed Top-K SS and a separate Top-1 retargeting PBR stage,
rather than requiring one aggregation operator to be best at both tasks.

The contributions are an explicit ray-shared source-adaptation formulation;
matched tests separating raw retrieval, decoder descriptors, scalar scores,
aggregation, and retargeting; and a stage-specific evaluation with a frozen
held-out protocol. Candidate coverage, component effects, and final pipeline
quality are separate claims. We retain negative results and uncertainty
instead of inferring necessity from the architecture alone.
